%% QDCVR — Knowledge-Based Systems submission (elsarticle)
%% Compile: pdflatex main && bibtex main && pdflatex main && pdflatex main
%% AUTHOR ACTION: replace author placeholders, ORCID, and funding fields.
\documentclass[preprint,12pt]{elsarticle}

\usepackage[T1]{fontenc}
\usepackage[utf8]{inputenc}
\usepackage{booktabs}
\usepackage{array}
\usepackage{graphicx}
\usepackage{amsmath}
\usepackage{amssymb}
\usepackage{xcolor}
\usepackage{url}
\usepackage[colorlinks=true,linkcolor=blue,citecolor=blue,urlcolor=blue]{hyperref}
\usepackage{caption}
\captionsetup{font=small,labelfont=bf}
\newcommand{\yes}{\checkmark}
\newcommand{\no}{--}

\graphicspath{{figs/}}
\setlength{\emergencystretch}{1.5em}

\journal{Knowledge-Based Systems}

\begin{document}

\begin{frontmatter}

\title{Design of a content-verified knowledge-base management system for
trustworthy enterprise document intelligence}

\author[inst1]{Author Name\corref{cor1}}
\ead{author@example.com}
\author[inst1]{Co-author Name}
\cortext[cor1]{Corresponding author.}
\affiliation[inst1]{organization={Affiliation Department, University,
City, Country}, city={City}, country={Country}}

\begin{abstract}
Enterprises accumulate process regulations, failure analyses, and operational
lessons as unstructured documents spread across shared drives. We present
QDCVR, a self-hosted knowledge-base management system that turns a folder of
PDFs, Office documents, and scans into knowledge bases that agents can be
trusted to answer from. The system spans six functional subsystems: a
quality-gated parsing and ingestion pipeline; knowledge-base management over a
five-layer data model; content-routed knowledge organization with tags and a
document graph; an operational-knowledge (experience) store with credibility
tiers and automated distillation; a multi-lane retrieval layer in which every
candidate citation must pass a content-verification gate; and an agent-facing
tool interface that exposes the whole platform as Model Context Protocol
tools. Its design principle is that organization, retrieval, and agency all
terminate at explicit contracts rather than implicit prompts, and that no
similarity score may promote a document into an answer. We evaluate the system
on a deployed industrial corpus (five knowledge bases, 42 documents, sixteen
gold-annotated questions) and through functional and non-functional test
campaigns. The verification gate exposes a same-domain score-inflation failure
consistent with known LLM-judge miscalibration: non-gold scores are
indistinguishable from gold (Mann--Whitney $p = 0.26$), and judged re-ranking
lowers Hit@5 from 0.94 to 0.44, while reusing the same verdicts as a filter
restores 0.94 (sign test $p = 0.008$). Functional campaigns confirm ingestion,
deduplication, organization, graph construction, experience distillation, and
abstention end to end; concurrency and stability campaigns show no data loss
under twenty-way concurrent writes. Verifier calibration is the open problem;
honesty needs an enforcement layer.
\end{abstract}

\begin{keyword}
Knowledge-base management system \sep Document ingestion \sep Knowledge
organization \sep Content-verified retrieval \sep Trustworthy question
answering \sep Model Context Protocol
\end{keyword}

\end{frontmatter}

%% ---------------------------------------------------------------- highlights
%% Submitted separately as highlights.tex; kept here for reference.

